\documentclass[11pt]{article}
\usepackage[margin=1in]{geometry}
\usepackage{amsmath,amssymb}
\usepackage{booktabs}
\usepackage{array}
\usepackage{hyperref}

\title{Non-Dimensional Analysis of the LuGre--Adamov Mecanum Plant\\
\large Running notes (updated as each equation is treated)}
\author{Vishvesh Koranne}
\date{\today}

\newcommand{\st}[1]{{#1}^{*}}

\begin{document}
\maketitle

\section*{Purpose}
Quantify the nonlinearity of the full mecanum plant (dynamic LuGre friction with
Adamov spin coupling, as implemented in \texttt{run\_one.jl},
\texttt{dynamics\_full\_mf\_asmc!} and \texttt{lugre\_dyn\_rates}) by varying
physically meaningful dimensionless constants, and from this determine which
classes of motion warrant nonlinear/robust control and for which a linear
controller suffices. Each equation is scaled by the physical quantity it
balances. See \texttt{intent.txt}.

\section{Notation and reference scales}
\begin{tabular}{@{}ll@{}}
$m$, $m_w$, $m_s = m + 4m_w$ & body mass, wheel mass, total mass \\
$I_s$ & platform moment of inertia about the vertical axis \\
$a_X$, $a_Y$ & centre-of-mass offset from the geometric centre \\
$h$, $l$ & half-length, half-width (wheel-centre positions $\pm h$, $\pm l$) \\
$\chi$ & roller contact-patch radius \\
$\mu$ & Coulomb friction coefficient \\
$V_0$, $\Omega_0$, $T$ & reference body speed, yaw rate, and manoeuvre time \\
\end{tabular}

\medskip
\textbf{Reference scales.} Force $m_s g$; length $h$ (half-length); moment
$m_s g h(1+\Lambda)$ for the yaw equation, where $h(1+\Lambda) = h + l$ is the yaw
lever of the mecanum O-configuration. The width enters only through the aspect
ratio $\Lambda = l/h$.

\medskip
Dimensionless states: $V_x = V_0\st{V_x}$, $V_y = V_0\st{V_y}$,
$\dot\psi = \Omega_0\st{\dot\psi}$, $t = T\st{t}$ (so
$\ddot\psi = (\Omega_0/T)\,\st{\ddot\psi}$, $\dot V = (V_0/T)\,\st{\dot V}$),
contact forces $F_{x,i} = m_s g\,\st{F_{x,i}}$, spin torques
$M_{z,i} = m_s g h(1+\Lambda)\,\st{M_{z,i}}$, wheel-centre positions
$p_{x,i} = s_{x,i}\,h$, $p_{y,i} = s_{y,i}\,\Lambda h$ with
$(s_{x,i}, s_{y,i}) \in \{(+,+),(+,-),(-,+),(-,-)\}$ for $i = 1..4$.

\subsection*{Dimensionless constants}
\textbf{Plant / environment} (vary with floor and payload placement):
\begin{equation}
\mu, \qquad
\varepsilon_X = \frac{a_X}{h}, \qquad
\varepsilon_Y = \frac{a_Y}{l} = \frac{a_Y}{\Lambda h}, \qquad
\hat\chi = \frac{\chi}{h}.
\end{equation}
\textbf{Fixed for a given robot} (change only with payload mass/geometry):
\begin{equation}
r_m = \frac{m}{m_s}, \qquad
\Lambda = \frac{l}{h}, \qquad
\iota = \frac{I_s}{m_s h^2(1+\Lambda)^2}, \qquad
\hat\jmath_w = \frac{J_w}{m_s R^2}, \qquad
\hat\tau = \frac{\tau_{\max}}{m_s g R},
\end{equation}
DC-motor constants (voltage-input plant, \texttt{hybrid\_ctrl/plant.jl}):
\begin{equation}
\hat\tau_V = \frac{G\eta K_t V_{\max}}{R_a^{\mathrm{m}}\,m_s g R}, \qquad
\hat\tau_i = \frac{G\eta K_t\, i_{\max}}{m_s g R}, \qquad
\hat\tau_f = \frac{G\eta\,\tau_f}{m_s g R}.
\end{equation}
\begin{equation}
\rho = \frac{R}{h}, \qquad
\rho_a = \frac{R_a}{R}, \qquad
\hat\jmath_r = \frac{n_r J_r}{m_s r_r^2},
\qquad r_r = R - R_a = \rho(1-\rho_a)\,h.
\end{equation}
Friction-model constants:
\begin{equation}
r_s = \frac{\mu_s}{\mu_c}, \qquad
\kappa_\sigma = \frac{\sigma_0}{\sigma_{0s}}, \qquad
\kappa_\chi = \frac{\chi}{r_r} = \frac{\hat\chi}{\rho(1-\rho_a)}\ \text{(derived)},
\end{equation}
\begin{equation}
\zeta_b = \frac{\sigma_1}{2}\sqrt{\frac{g}{\sigma_0}}, \qquad
\zeta_{bs} = \frac{\sigma_{1s}}{2}\sqrt{\frac{g}{\sigma_{0s}}}
\qquad \text{(pre-sliding damping ratios)}.
\end{equation}
Derived lever fractions (functions of $\Lambda$ only, not new variables):
$\lambda_h = 1/(1+\Lambda)$, $\lambda_l = \Lambda/(1+\Lambda)$, $\lambda_h + \lambda_l = 1$.
\textbf{Motion (grip demand)} (vary with the commanded trajectory):
\begin{equation}
\Pi_a = \frac{V_0}{\mu g T},\qquad
\Pi_c = \frac{\Omega_0 V_0}{\mu g},\qquad
\Pi_\omega = \frac{h\Omega_0^2}{\mu g},\qquad
\Pi_\alpha = \frac{h\Omega_0}{\mu g T},
\end{equation}
\begin{equation}
\beta_1 = \frac{p_1 V_0}{\mu m_s g R^2},\qquad
\beta_m = \frac{G^2\eta K_t K_b\,V_0}{R_a^{\mathrm{m}}\,\mu m_s g R^2},\qquad
\beta_2 = \frac{p_2 V_0}{\mu m_s g r_r^2},
\end{equation}
\begin{equation}
\Pi_b = \frac{\mu}{\sigma_0 V_0 T},\qquad
\hat v = \frac{v_{\mathrm{str}}}{V_0},\qquad
\hat\omega_b = T\sqrt{\sigma_0 g},
\end{equation}
and the derived motion ratio $\nu = h\Omega_0/V_0 = \Pi_\omega/\Pi_c$
(rotational vs.\ translational speed at the wheel; not a new variable).

\medskip
\begin{tabular}{@{}llll@{}}
\toprule
Constant & Base value & Constant & Base value \\
\midrule
$r_m$ & 0.843 & $\varepsilon_X$ & 0.068 \\
$\Lambda$ & 0.638 & $\varepsilon_Y$ & $-0.173$ \\
$\iota$ & 0.838 & $\hat\chi$ & 0.021 \\
$\hat\jmath_w$ & 0.066 & $\mu$ & $\{0.3, 0.5, 0.8\}$ \\
$\hat\tau$ & 0.573 & $\rho$ & 0.213 \\
$\hat\jmath_r$ & $1.60\times10^{-3}$ & $\rho_a$ & 0.710 \\
$r_s$ & 1.1 & $\kappa_\sigma$ & 1.505 \\
$\kappa_\chi$ & 0.345 & $\zeta_b$ & 0.062 \\
$\zeta_{bs}$ & 0.052 & $\hat\tau_V$ & 0.535 \\
$\hat\tau_i$ & 0.571 & $\hat\tau_f$ & 0.013 \\
\bottomrule
\end{tabular}

\section{Body translational equations ($V_x$, $V_y$)}

\subsection{Dimensional form}
Force balance along the body axes (N):
\begin{align}
m_s\dot V_x - m a_Y\ddot\psi &= \textstyle\sum_i F_{x,i}
  + m_s\dot\psi V_y + m a_X\dot\psi^2, \label{eq:x}\\
m_s\dot V_y + m a_X\ddot\psi &= \textstyle\sum_i F_{y,i}
  - m_s\dot\psi V_x + m a_Y\dot\psi^2. \label{eq:y}
\end{align}
The only external force on the body in the plane is the contact traction
$\sum_i F_i$, so the balanced quantity is traction force.

\subsection{Scaling choice}
\begin{itemize}
  \item \textbf{Force scale: $m_s g$} (total weight $= \sum_i N_i$, since the
    normal loads include the wheel weights). The friction force is kept on this
    scale so that $\mu$ remains a free parameter that can be varied across
    simulations; $\sum_i \st{F_i}$ is bounded by roughly $\mu_s$.
  \item \textbf{Inertial terms are multiplied and divided by $\mu$}, so that
    each one reads as $\mu\times$(fraction of available grip demanded by the
    motion).
\end{itemize}

\subsection{$\mu$-modified non-dimensional form}
Dividing \eqref{eq:x}--\eqref{eq:y} by $m_s g$:
\begin{align}
\mu\Big[\Pi_a\,\st{\dot V_x} - r_m\Lambda\varepsilon_Y\Pi_\alpha\,\st{\ddot\psi}\Big]
 &= \textstyle\sum_i \st{F_{x,i}}
 + \mu\Big[\Pi_c\,\st{\dot\psi}\st{V_y} + r_m\varepsilon_X\Pi_\omega\,(\st{\dot\psi})^{2}\Big],
 \label{eq:xnd}\\
\mu\Big[\Pi_a\,\st{\dot V_y} + r_m\varepsilon_X\Pi_\alpha\,\st{\ddot\psi}\Big]
 &= \textstyle\sum_i \st{F_{y,i}}
 + \mu\Big[-\Pi_c\,\st{\dot\psi}\st{V_x} + r_m\Lambda\varepsilon_Y\Pi_\omega\,(\st{\dot\psi})^{2}\Big].
 \label{eq:ynd}
\end{align}

\subsection{Significance of each term}
\begin{tabular}{@{}>{$}l<{$}p{0.66\linewidth}@{}}
\toprule
\text{Term} & Physical significance \\
\midrule
\mu\,\Pi_a\,\st{\dot V}
  & Linear inertia: acceleration $V_0/T$ as a fraction $\Pi_a$ of the grip-limited
    acceleration $\mu g$. \\
\mu\,r_m(\Lambda\varepsilon_Y,\ \varepsilon_X)\,\Pi_\alpha\,\st{\ddot\psi}
  & Yaw-acceleration coupling through the CoM offset: tangential acceleration of
    the offset CoM vs.\ $\mu g$; vanishes for a centred CoM. \\
\sum_i \st{F_i}
  & Net contact traction vs.\ total weight; bounded by $\approx\mu_s$. Carries all
    of the friction nonlinearity (pre-sliding $\to$ Stribeck $\to$ gross sliding). \\
\mu\,\Pi_c\,\st{\dot\psi}\st{V}
  & Coriolis/centripetal demand of turning while translating, $\Omega_0V_0$, vs.\
    $\mu g$. Bilinear in the states: a geometric (non-friction) nonlinearity. \\
\mu\,r_m(\varepsilon_X,\ \Lambda\varepsilon_Y)\,\Pi_\omega\,(\st{\dot\psi})^{2}
  & Centripetal load on the offset CoM during rotation vs.\ $\mu g$; quadratic in
    yaw rate; vanishes for a centred CoM. \\
\bottomrule
\end{tabular}

\medskip
When the bracketed grip-demand groups are $\ll 1$ the contacts stay in the stiff,
near-linear pre-sliding regime; as the combined demand approaches $1$ the contacts
enter the Stribeck and gross-sliding regimes and friction saturates.

\subsection{Normal-load distribution}
The CoM offset also sets the static normal loads, so $\varepsilon_X$,
$\varepsilon_Y$ enter the right-hand side through the friction forces:
\begin{equation}
\st{N_i} = \frac{N_i}{m_s g}
 = \frac{r_m}{4}\big(1 + s_{x,i}\,\varepsilon_X + s_{y,i}\,\varepsilon_Y\big)
 + \frac{1-r_m}{4}.
\end{equation}
This decides which wheel saturates first and may matter more for nonlinearity
than the small CoM terms on the left-hand side.

\section{Body yaw equation ($\dot\psi$)}

\subsection{Dimensional form}
Moment balance about the vertical axis through the geometric centre (N\,m):
\begin{equation}
I_s\ddot\psi - m a_Y\dot V_x + m a_X\dot V_y
 = \textstyle\sum_i \big(p_{x,i}F_{y,i} - p_{y,i}F_{x,i}\big)
 + \sum_i M_{z,i}
 - m\dot\psi\,(a_X V_x + a_Y V_y). \label{eq:yaw}
\end{equation}

\subsection{Scaling choice}
Moment scale $m_s g h(1+\Lambda)$: the force scale of the translational equations
times the yaw lever $h + l$, so that both planar dimensions enter. With
$45^\circ$ rollers each wheel's traction acts on the yaw axis through this lever,
so the available yaw moment scales as $\mu\,m_s g h(1+\Lambda)$. Inertial terms
are again multiplied and divided by $\mu$.

\subsection{$\mu$-modified non-dimensional form}
\begin{multline}
\mu\Big[\iota(1+\Lambda)\,\Pi_\alpha\,\st{\ddot\psi}
 - r_m\lambda_l\varepsilon_Y\Pi_a\,\st{\dot V_x}
 + r_m\lambda_h\varepsilon_X\Pi_a\,\st{\dot V_y}\Big]
 = \textstyle\sum_i\big(\lambda_h s_{x,i}\st{F_{y,i}} - \lambda_l s_{y,i}\st{F_{x,i}}\big)
 + \sum_i \st{M_{z,i}} \\
 - \mu\,r_m\Pi_c\,\st{\dot\psi}\big(\lambda_h\varepsilon_X\st{V_x} + \lambda_l\varepsilon_Y\st{V_y}\big).
 \label{eq:yawnd}
\end{multline}

\subsection{Significance of each term}
\begin{tabular}{@{}>{$}l<{$}p{0.62\linewidth}@{}}
\toprule
\text{Term} & Physical significance \\
\midrule
\mu\,\iota(1+\Lambda)\,\Pi_\alpha\,\st{\ddot\psi}
  & Yaw inertia: tangential acceleration $(h+l)\Omega_0/T$ at the yaw lever vs.\
    $\mu g$, weighted by $\iota = (k/(h+l))^2$, the squared ratio of the
    platform's radius of gyration $k$ to the yaw lever. \\
\mu\,r_m(\lambda_l\varepsilon_Y,\ \lambda_h\varepsilon_X)\,\Pi_a\,\st{\dot V}
  & Linear acceleration of the offset CoM producing a yaw moment; vanishes for a
    centred CoM. \\
\sum_i(\lambda_h s_x\st{F_y} - \lambda_l s_y\st{F_x})
  & Yaw moment of the contact traction forces at the wheel positions; carries the
    translational friction nonlinearity into yaw. $\lambda_h$, $\lambda_l$ are
    the shares of the yaw lever from the length and the width. \\
\sum_i \st{M_{z,i}}
  & Contact-patch spin friction torque. Bounded by
    $\approx(3\pi/16)\,\mu_s\,\lambda_h\hat\chi \approx 0.008\,\mu$ in total: small, but the
    only yaw term that comes purely from spin friction. \\
\mu\,r_m\Pi_c\,\st{\dot\psi}(\lambda_h\varepsilon_X\st{V_x} + \lambda_l\varepsilon_Y\st{V_y})
  & Gyroscopic term from the CoM offset (turning while translating); bilinear in
    the states; vanishes for a centred CoM. \\
\bottomrule
\end{tabular}

\subsection{Yaw inertia and the COM offset}
$I_s$ is the inertia about the geometric centre, the reference point of the mass
matrix. When the COM offset changes because a payload is moved, it changes with
it by the parallel-axis theorem:
\begin{equation}
I_s = I_G + m(a_X^2 + a_Y^2), \qquad
\iota = \iota_G + r_m\,\frac{\varepsilon_X^2 + \Lambda^2\varepsilon_Y^2}{(1+\Lambda)^2},
\qquad \iota_G = \frac{I_G}{m_s h^2(1+\Lambda)^2},
\end{equation}
so $\varepsilon_X$, $\varepsilon_Y$ and $\iota$ are not independent: $\iota_G$ is
the fixed robot constant and $\iota$ follows from the offset.

\section{Wheel spin equation ($\omega_i$)}

\subsection{Dimensional form}
Torque balance about the wheel axle (N\,m), one per wheel:
\begin{equation}
J_w\dot\omega_i = \tau_{\mathrm{sat},i} - F_{x,i}R - p_1\omega_i,
\qquad \tau_{\mathrm{sat},i} = \mathrm{sat}(\tau_i, \tau_{\max}). \label{eq:wheel}
\end{equation}

\subsection{Scaling choice}
Torque scale $m_s g R$: the force scale of the body equations times the wheel
radius, so the rim reaction $F_{x,i}R$ becomes exactly the same $\st{F_{x,i}}$
that appears in the body equations. Wheel speed is scaled by the rolling speed,
$\omega_i = (V_0/R)\,\st{\omega_i}$. The inertial and viscous terms are
multiplied and divided by $\mu$.

\subsection{$\mu$-modified non-dimensional form}
\begin{equation}
\mu\,\hat\jmath_w\Pi_a\,\st{\dot\omega_i}
 = \st{\tau_{\mathrm{sat},i}} - \st{F_{x,i}} - \mu\,\beta_1\,\st{\omega_i},
\qquad |\st{\tau_{\mathrm{sat},i}}| \le \hat\tau. \label{eq:wheelnd}
\end{equation}

\subsection{Significance of each term}
\begin{tabular}{@{}>{$}l<{$}p{0.66\linewidth}@{}}
\toprule
\text{Term} & Physical significance \\
\midrule
\mu\,\hat\jmath_w\Pi_a\,\st{\dot\omega}
  & Wheel inertia: $\hat\jmath_w$ is the wheel's rotational inertia expressed as an
    equivalent translating mass, as a fraction of the total mass. \\
\st{\tau_{\mathrm{sat}}}
  & Motor torque vs.\ weight torque; saturates at $\hat\tau$ (actuator
    nonlinearity). \\
\st{F_x}
  & Contact friction reaction at the rim, bounded by $\approx\mu_s\st{N_i}$
    (friction nonlinearity). \\
\mu\,\beta_1\,\st{\omega}
  & Drivetrain viscous torque at the reference speed as a fraction of the grip
    torque $\mu m_s g R$. \\
\bottomrule
\end{tabular}

\subsection{Two saturations}
Per wheel the motor can deliver up to $\hat\tau = 0.573$, while the contact can
transmit only $\approx\mu_s\st{N_i}\approx\mu/4$. Their ratio
$\hat\tau/(\mu/4)\approx 2.3/\mu$ is $2.9$--$7.6$ for $\mu = 0.8$--$0.3$, so the
motor can overpower the contact and drive the wheel into slip well before the
actuator itself saturates. The wheel equation behaves linearly only while both
$\st{\tau}$ and $\st{F_x}$ stay inside their limits.

\subsection{DC-motor torque map (back-EMF limit)}
In the voltage-input plant (\texttt{hybrid\_ctrl/plant.jl}, \texttt{motor\_torque})
$\tau_{\mathrm{sat},i}$ is not a fixed cap but the output of a static motor map.
With armature inductance neglected ($L_a = 0$), the armature voltage balance
(V) and the gearbox torque relation are
\begin{equation}
V_i = R_a^{\mathrm{m}}\,i_i + K_b\,G\,\omega_i,\qquad
|i_i| \le i_{\max},\quad |V_i| \le V_{\max},\qquad
\tau_i = G\eta\big(K_t\,i_i - \tau_f\,\mathrm{sgn}(G\omega_i)\big),
\end{equation}
where $G$ is the gear ratio, $\eta$ the gear efficiency, $K_t$, $K_b$ the torque
and back-EMF constants, $R_a^{\mathrm{m}}$ the winding resistance (superscript to
distinguish it from the roller axle distance $R_a$) and $\tau_f$ the motor-shaft
Coulomb friction. A voltage-rate (slew) limit $|\dot V| \le \dot V_{\max}$ is
applied in the mixer.

\paragraph{Non-dimensional form.} With $V_i = V_{\max}\st{V_i}$ and the wheel
torque scale $m_s g R$:
\begin{equation}
\st{\tau_i} = \mathrm{clamp}\Big(\hat\tau_V\,\st{V_i} - \mu\,\beta_m\,\st{\omega_i},\
 -\hat\tau_i,\ \hat\tau_i\Big) - \hat\tau_f\,\mathrm{sgn}(\st{\omega_i}),
\qquad |\st{V_i}| \le 1 .
\end{equation}

\medskip
\begin{tabular}{@{}>{$}l<{$}p{0.70\linewidth}@{}}
\toprule
\text{Constant} & Physical significance \\
\midrule
\hat\tau_V = 0.535
  & Stall torque available from the full bus voltage, vs.\ weight torque. \\
\hat\tau_i = 0.571
  & Current-limited torque vs.\ weight torque. $\hat\tau_V < \hat\tau_i$, so with
    the base $R_a^{\mathrm{m}} = 2\,\Omega$ the voltage limit binds first. \\
\mu\beta_m\st{\omega}
  & Back-EMF torque loss at the reference speed vs.\ the grip torque. The
    electrical damping $G^2\eta K_tK_b/R_a^{\mathrm{m}} = 0.337$\,N\,m\,s is about
    three times the drivetrain viscous $p_1 = 0.11$. \\
\hat\tau_f = 0.013
  & Motor-shaft Coulomb friction vs.\ weight torque (sign nonlinearity at
    wheel reversals). \\
V_{\max}/(\dot V_{\max}T)
  & Voltage slew time ($0.12$\,s) vs.\ manoeuvre time. \\
\bottomrule
\end{tabular}

\paragraph{Speed-dependent torque envelope.} The torque available at full
voltage falls linearly with wheel speed,
\begin{equation}
|\st{\tau}|_{\mathrm{avail}} = \min\big(\hat\tau_i,\ \hat\tau_V - \mu\beta_m|\st\omega|\big)
 = \hat\tau_V\Big(1 - \frac{R|\omega|}{v_{\mathrm{nl}}}\Big),
\qquad v_{\mathrm{nl}} = \frac{R\,V_{\max}}{K_b G} = 1.39\ \text{m/s},
\end{equation}
reaching zero at the no-load rim speed $v_{\mathrm{nl}}$ (in non-dimensional form
$v_{\mathrm{nl}}/V_0 = \hat\tau_V/(\mu\beta_m)$, not a new variable). Comparing
with the per-wheel grip $\approx\mu\st{N_i}\approx\mu/4$, the motor envelope drops
below the grip limit for rim speeds above
\begin{equation}
\frac{R|\omega|}{v_{\mathrm{nl}}} \gtrsim 1 - \frac{\mu}{4\hat\tau_V}
 \approx 1 - 0.47\,\mu ,
\end{equation}
i.e.\ about $1.06$\,m/s at $\mu = 0.5$ and $0.87$\,m/s at $\mu = 0.8$. Below this
speed friction saturates first (wheel slip); above it the actuator saturates
first. The two saturations of the wheel equation therefore trade places with
speed, and the crossover depends only on $\mu$ and $\hat\tau_V$.

\paragraph{Torque-input plant.} \texttt{run\_one.jl} bypasses this map and applies
a smooth saturation at $\tau_{\max} = 10$\,N\,m ($\hat\tau = 0.573$) with no
back-EMF; it corresponds to the limit $\beta_m \to 0$, $\hat\tau_V \to \hat\tau$.

\subsection{Wheel inertia reflected into the body}
Under rolling without slip, the four wheels add their inertia to the body
equations: the effective mass becomes $m_s(1 + 4\hat\jmath_w)$ in $x$ and $y$,
and the yaw inertia constant becomes
\begin{equation}
\iota_{\mathrm{aug}} = \iota + 4\hat\jmath_w = 0.838 + 0.264 = 1.102,
\end{equation}
consistent with \texttt{M\_aug} in \texttt{run\_one.jl}
($m_{\tilde{}} = m_s + 4J_w/R^2$, $I_{\psi,\mathrm{aug}} = I_s + 4J_w(h+l)^2/R^2$).
With the yaw lever $h+l$ as the reference, the same $4\hat\jmath_w$ appears in
both translation and yaw.

\section{Roller spin equation ($\dot\gamma_i$)}

\subsection{Dimensional form}
Torque balance about the axle of the roller in contact (N\,m), one per wheel,
with the inertia of the $n_r$ rollers lumped:
\begin{equation}
n_r J_r\ddot\gamma_i = -p_2\dot\gamma_i - F_{x,i}A_i - F_{y,i}B_i - M_{z,i}C_i,
\label{eq:roller}
\end{equation}
\begin{align*}
A_i &= \sin\delta_i\,(R - R_a\cos\tilde\theta_i) - DY_i\sin\tilde\theta_i\cos\delta_i,
\qquad DY_i = R_a\tan\delta_i\tan\tilde\theta_i,\\
B_i &= \cos\delta_i\,(R_a - R\cos\tilde\theta_i), \qquad
C_i = \sin\tilde\theta_i\cos\delta_i,
\end{align*}
where $\tilde\theta_i = \mathrm{saw}(\theta_i)\in[-15^\circ, 15^\circ]$ locates the
contact on the roller in contact.

\subsection{Scaling choice}
Torque scale $m_s g\,r_r$: the force scale times the roller radius
$r_r = R - R_a$, which is the lever of the contact forces about the roller axle.
Roller spin is scaled by the speed at which the roller surface must move to
carry the reference speed, $\dot\gamma_i = (V_0/r_r)\,\st{\dot\gamma_i}$.
Inertial and viscous terms are multiplied and divided by $\mu$.

\subsection{Dimensionless lever arms}
Using $\tan\delta\cos\delta = \sin\delta$, the levers simplify to
\begin{equation}
\hat A_i = \frac{A_i}{r_r} = \sin\delta_i\,\frac{1 - \rho_a/\cos\tilde\theta_i}{1-\rho_a},
\qquad
\hat B_i = \frac{B_i}{r_r} = -\cos\delta_i\,\frac{\cos\tilde\theta_i - \rho_a}{1-\rho_a},
\qquad
\hat C_i = \sin\tilde\theta_i\cos\delta_i .
\end{equation}
At the roller centre ($\tilde\theta = 0$): $\hat A = \sin\delta$,
$\hat B = -\cos\delta$, $\hat C = 0$. At the roller edge ($\tilde\theta = \pm15^\circ$):
$\hat A = 0.91\sin\delta$, $\hat B = -0.88\cos\delta$,
$\hat C = \pm0.26\cos\delta$. The levers vary periodically as each roller passes
through contact, at the roller-passing frequency $n_r\omega_i/2\pi$.

\subsection{$\mu$-modified non-dimensional form}
\begin{equation}
\mu\,\hat\jmath_r\Pi_a\,\st{\ddot\gamma_i}
 = -\mu\,\beta_2\,\st{\dot\gamma_i}
   - \hat A_i\,\st{F_{x,i}} - \hat B_i\,\st{F_{y,i}}
   - \kappa_s\,\hat C_i\,\st{M_{z,i}},
\qquad \kappa_s = \frac{h(1+\Lambda)}{r_r} = \frac{1+\Lambda}{\rho(1-\rho_a)}.
\label{eq:rollernd}
\end{equation}
($\kappa_s = 26.5$ only converts $\st{M_z}$ from the yaw scale $m_s g h(1+\Lambda)$
to the roller scale; the spin torque itself is bounded by
$\approx(3\pi/16)\mu_s N_i\chi$, i.e.\ $\approx 0.6\,\mu_s\st{N_i}\,\chi/r_r$ on the
roller scale, with $\chi/r_r = 0.34$.)

\subsection{Significance of each term}
\begin{tabular}{@{}>{$}l<{$}p{0.62\linewidth}@{}}
\toprule
\text{Term} & Physical significance \\
\midrule
\mu\,\hat\jmath_r\Pi_a\,\st{\ddot\gamma}
  & Roller inertia, as an equivalent translating mass at the roller surface, as a
    fraction of the total mass ($\hat\jmath_r = 1.6\times10^{-3}$, about 40 times
    smaller than $\hat\jmath_w$). \\
\mu\,\beta_2\,\st{\dot\gamma}
  & Roller-bearing viscous torque at the reference roller speed as a fraction of
    the grip torque $\mu m_s g r_r$. Measures how far the roller is from an ideal
    free roller. \\
\hat A\st{F_x} + \hat B\st{F_y}
  & Torque of the contact friction forces about the roller axle. The lever arms
    depend on where the contact sits on the roller (geometric, periodic
    nonlinearity through $\tilde\theta$). \\
\kappa_s\hat C\,\st{M_z}
  & Spin friction torque projected onto the roller axle; zero at the roller
    centre, largest at the roller edges. \\
\bottomrule
\end{tabular}

\subsection{Quasi-static (slaved) roller}
The ratio of the inertial to the viscous term is
$\hat\jmath_r\Pi_a/\beta_2 = \tau_r/T$, with the roller mechanical time constant
$\tau_r = n_r J_r/p_2 \approx 2.1$\,ms. For any manoeuvre with $T \gg \tau_r$ the
roller equation is effectively algebraic,
\begin{equation}
\mu\,\beta_2\,\st{\dot\gamma_i}
 \approx -\hat A_i\st{F_{x,i}} - \hat B_i\st{F_{y,i}} - \kappa_s\hat C_i\st{M_{z,i}},
\end{equation}
i.e.\ the roller spins until the friction torque about its axle balances bearing
drag. Since $\beta_2$ is small, the friction force along the roller's rolling
direction stays small: this is the free-roller constraint that gives the
mecanum wheel its force direction. $\tau_r/T$ is the small parameter of a
singular perturbation, and the roller state is a source of stiffness rather than
of slow dynamics.

\section{Bristle equations ($z_x$, $z_y$, $z_s$)}

\subsection{Dimensional form}
Rate balance of the bristle deflection at each contact: slip drives the
deflection, relaxation pulls it back (translational in m/s, spin in rad/s):
\begin{align}
\dot z_{x} &= V_{px} - \sigma_0\,\frac{s_t}{g_t}\,z_{x}, \qquad
\dot z_{y} = V_{py} - \sigma_0\,\frac{s_t}{g_t}\,z_{y}, \label{eq:bristle}\\
\dot z_{s} &= \omega_{z} - \sigma_{0s}\,\frac{s_s}{g_s}\,z_{s}, \label{eq:bristles}
\end{align}
with effective sliding speeds (Adamov coupling) and Stribeck curves
\begin{gather*}
s_t = f_{sl}\,\tfrac{8}{3\pi}\chi|\omega_z| + |V_p|, \qquad
s_s = \tfrac{16}{3\pi}\chi|\omega_z| + 5|V_p|,\\
g_{t,s} = \mu\big[1 + (r_s - 1)\,e^{-(s/v_{\mathrm{str}})^2}\big], \qquad
f_{sl} = 1 - \big(1 - |z|/\delta^*\big)^{2/3},\quad \delta^* = \mu_s/\sigma_0 .
\end{gather*}

\subsection{Scaling choice}
\begin{itemize}
  \item \textbf{Deflection} relative to the Coulomb breakaway deflection
    $\delta_c = \mu/\sigma_0$:
    $\hat z = \sigma_0 z/\mu$. Since the elastic friction force is
    $N\sigma_0 z$, $\mu\hat z$ is the instantaneous elastic friction
    coefficient, and $|\hat z| \lesssim r_s$.
  \item \textbf{Spin deflection} relative to its breakaway in the same way:
    $\hat z_s = \chi\sigma_{0s} z_s/\mu$, so that $\mu\hat z_s = M_z^{\mathrm{el}}/(N\chi)$.
  \item \textbf{Slip and sliding speeds} by $V_0$; \textbf{contact spin} by the
    roller-surface rate, $\omega_z = (V_0/r_r)\,\st{\omega_z}$.
  \item \textbf{Stribeck curves} are written as $g = \mu\,\hat g$ with the
    $\mu$-free shape $\hat g(\st s) = 1 + (r_s-1)\,e^{-(\st s/\hat v)^2}$.
\end{itemize}

\subsection{Non-dimensional form}
\begin{align}
\Pi_b\,\hat z_{x}' &= \st{V_{px}} - \frac{\st{s_t}}{\hat g_t}\,\hat z_{x}, \qquad
\Pi_b\,\hat z_{y}' = \st{V_{py}} - \frac{\st{s_t}}{\hat g_t}\,\hat z_{y},
\label{eq:bristlend}\\
\kappa_\sigma\Pi_b\,\hat z_{s}' &= \kappa_\chi\,\st{\omega_z}
 - \frac{\st{s_s}}{\hat g_s}\,\hat z_{s}, \label{eq:bristlesnd}
\end{align}
\begin{gather*}
\st{s_t} = f_{sl}\,\tfrac{8}{3\pi}\kappa_\chi|\st{\omega_z}| + |\st{V_p}|, \qquad
\st{s_s} = \tfrac{16}{3\pi}\kappa_\chi|\st{\omega_z}| + 5|\st{V_p}|, \qquad
f_{sl} = 1 - \big(1 - |\hat z|/r_s\big)^{2/3}.
\end{gather*}
Here $(\cdot)'$ is the derivative with respect to $\st t = t/T$. $\mu$ has
dropped out of the equations except through $\Pi_b$.

\subsection{Significance of each term and constant}
\begin{tabular}{@{}>{$}l<{$}p{0.66\linewidth}@{}}
\toprule
\text{Term / constant} & Physical significance \\
\midrule
\Pi_b\,\hat z'
  & Bristle storage. $\Pi_b = \delta_c/(V_0T)$ is the pre-sliding (breakaway)
    length compared with the distance travelled in a manoeuvre; equivalently the
    bristle relaxation time $\delta_c/V_0$ vs.\ $T$. \\
\st{V_p}
  & Contact slip velocity driving the deflection. \\
\st{s_t}\hat z/\hat g_t
  & Bristle relaxation. Bilinear in slip speed and deflection: the source of the
    pre-sliding $\to$ sliding transition. \\
\kappa_\chi\st{\omega_z}
  & Contact spin as a patch-edge speed $\chi|\omega_z|$ compared with $V_0$
    ($\kappa_\chi = \chi/r_r$). \\
\kappa_\sigma
  & Ratio of translational to spin bristle stiffness; the spin bristles relax
    $\kappa_\sigma$ times more slowly. \\
r_s,\ \hat v
  & Stiction-to-Coulomb ratio and Stribeck speed relative to $V_0$: the size and
    the width of the friction peak at low sliding speed. \\
\kappa_\chi|\st{\omega_z}|/|\st{V_p}|
  & Adamov coupling: when contact spin is large compared with translational slip,
    $\st{s_t}$ is raised and translational friction is reduced (spin
    ``lubricates'' translation). \\
\bottomrule
\end{tabular}

\subsection{Where the bristle equation is linear}
With $\Pi_b \approx 6\times10^{-4}$ ($\mu = 0.5$, $V_0 = 0.5$\,m/s, $T = 1$\,s)
the bristle states are fast. The equations are linear (a stiff spring,
$\hat z' \approx \st{V_p}/\Pi_b$, force $\propto$ deflection) only while the
relaxation term is negligible, i.e.\ while $|\hat z| \ll 1$: the slip
displacement accumulated over the manoeuvre must stay well below the breakaway
length $\delta_c = \mu/\sigma_0 \approx 0.2$--$0.5$\,mm. In terms of slip
velocity, $|V_p| \lesssim \delta_c/T$, a fraction of a mm/s. Once slip exceeds
this, relaxation dominates and the force saturates toward
$\mu N\,\hat g\,V_p/s$: Coulomb-like, sign-type behaviour. The numerical
regularisation $\varepsilon_{\mathrm{reg}} = 10^{-4}$\,m/s is of the same order
as this threshold, so part of the apparent ``linear zone'' is set by the
regularisation rather than by physics.

\section{Friction force laws}

\subsection{Dimensional form}
Force balance at the contact patch (N) and torque balance about the contact
normal (N\,m), per wheel:
\begin{gather}
F_{x} = -N\big(\sigma_0 z_x + \sigma_1\dot z_x + \sigma_2 V_{px}\big),\qquad
F_{y} = -N\big(\sigma_0 z_y + \sigma_1\dot z_y + \sigma_2 V_{py}\big),\\
M_z = -N\chi^2\big(\sigma_{0s}z_s + \sigma_{1s}\dot z_s + \sigma_{2s}\omega_z\big).
\end{gather}
In \texttt{run\_one.jl} $\sigma_2 = \sigma_{2s} = 0$ (viscous losses are carried by
$p_1$, $p_2$).

\subsection{Non-dimensional form}
Using the bristle scalings of the previous section,
$\dot z = V_0(\st{V_p} - \st{s_t}\hat z/\hat g_t)$ and
$\chi\dot z_s = V_0(\kappa_\chi\st{\omega_z} - \st{s_s}\hat z_s/\hat g_s)$, so
\begin{align}
\st{F_{x}} &= -\mu\,\st{N}\Big[\hat z_x
  + \Pi_{\sigma}\big(\st{V_{px}} - \tfrac{\st{s_t}}{\hat g_t}\hat z_x\big)\Big],
  \qquad \text{(same for $y$)}\\
\st{M_z} &= -\mu\,\st{N}\,\lambda_h\hat\chi\Big[\hat z_s
  + \Pi_{\sigma s}\big(\kappa_\chi\st{\omega_z} - \tfrac{\st{s_s}}{\hat g_s}\hat z_s\big)\Big],
\end{align}
with the micro-damping groups
\begin{equation}
\Pi_\sigma = \frac{\sigma_1 V_0}{\mu} = \frac{2\zeta_b}{\Pi_b\,\hat\omega_b},
\qquad
\Pi_{\sigma s} = \frac{\sigma_{1s} V_0}{\mu}
 = \frac{2\zeta_{bs}}{\sqrt{\kappa_\sigma}\,\Pi_b\,\hat\omega_b}.
\end{equation}

\subsection{Significance}
\begin{tabular}{@{}>{$}l<{$}p{0.66\linewidth}@{}}
\toprule
\text{Term / constant} & Physical significance \\
\midrule
\mu\st{N}\hat z
  & Elastic bristle force: the friction force carried by contact deflection. \\
\mu\st{N}\Pi_\sigma\,\dot{\hat z}\text{-term}
  & Bristle micro-damping. Vanishes in steady sliding ($\dot z \to 0$); active
    only during pre-sliding and transients. \\
\hat\omega_b = T\sqrt{\sigma_0 g}
  & Manoeuvre time in units of the pre-sliding oscillation period: the body
    mass bouncing on the bristle stiffness $N\sigma_0$
    ($\sqrt{\sigma_0 g} \approx 127$\,rad/s, about 20\,Hz). \\
\zeta_b,\ \zeta_{bs}
  & Damping ratio of that pre-sliding oscillation (translation, spin): a material
    property, independent of $\mu$ and of the motion. \\
\mu\st{N}\lambda_h\hat\chi\,\hat z_s
  & Spin friction torque relative to the yaw scale; contact-patch radius
    $\chi$ as the lever. \\
\bottomrule
\end{tabular}

\medskip
The pre-sliding regime is therefore a lightly damped ($\zeta_b \approx 0.06$)
oscillator at about 20\,Hz. In this regime the friction is linear but resonant;
outside it the friction saturates. Both are relevant to controller choice: a
linear design must either stay out of the saturation or tolerate the lightly
damped stick mode.

\section{Contact-slip kinematics}

\subsection{Dimensional form}
Velocity compatibility at the contact point: body velocity at the contact,
minus wheel rim speed, minus roller surface speed (m/s, rad/s), per wheel:
\begin{align*}
V_{px} &= V_x - \dot\psi(p_y + DY) - \omega R
  + \dot\gamma\sin\delta(R_a\cos\tilde\theta - R) + DY\dot\gamma\cos\delta\sin\tilde\theta,\\
V_{py} &= V_y + \dot\psi\,p_x + \dot\gamma\cos\delta(R\cos\tilde\theta - R_a),\\
\omega_z &= \dot\psi + \dot\gamma\sin\tilde\theta\cos\delta .
\end{align*}

\subsection{Non-dimensional form}
With $V = V_0\st V$, $\dot\psi = \Omega_0\st{\dot\psi}$,
$\omega = (V_0/R)\st\omega$, $\dot\gamma = (V_0/r_r)\st{\dot\gamma}$,
$\omega_z = (V_0/r_r)\st{\omega_z}$:
\begin{align}
\st{V_{px}} &= \st{V_x} - \nu\big(\Lambda s_y + \rho\rho_a\tan\delta\tan\tilde\theta\big)\st{\dot\psi}
  - \st{\omega} - \hat A\,\st{\dot\gamma},\\
\st{V_{py}} &= \st{V_y} + \nu\,s_x\,\st{\dot\psi} - \hat B\,\st{\dot\gamma},\\
\st{\omega_z} &= \nu\rho(1-\rho_a)\,\st{\dot\psi} + \hat C\,\st{\dot\gamma}.
\end{align}
The roller levers $\hat A$, $\hat B$, $\hat C$ are exactly those of the roller
equation.

\subsection{Significance}
\begin{itemize}
  \item The slip velocities are \emph{linear} in the velocity states; the only
    nonlinearity is the periodic dependence of the coefficients on
    $\tilde\theta = \mathrm{saw}(\theta)$ (a linear periodically time-varying
    structure at the roller-passing frequency).
  \item $\nu = h\Omega_0/V_0$ is the only motion group here: how much of the
    contact velocity comes from rotation vs.\ translation. Pure translation
    $\nu = 0$; pure spin-in-place $\nu \to \infty$. It is the natural axis for
    classifying motions.
  \item Virtual-power check: the $\dot\gamma$ coefficients in $\st{V_{px}}$,
    $\st{V_{py}}$ are $-\hat A$, $-\hat B$, matching the roller equation's
    $-\hat A\st{F_x} - \hat B\st{F_y}$. The $\dot\gamma$ coefficient in
    $\st{\omega_z}$ is $+\hat C$, so energy consistency requires $+\hat C M_z$ in
    the roller equation, whereas the code has $-\hat C M_z$ (see open notes).
\end{itemize}

\section{Material comparison: polyurethane vs.\ rubber rollers}
The friction constants are fixed for a given roller material, but comparing two
materials shows which ratios a material change would move. The rubber column is
an \emph{estimate} scaled from the PU baseline with Hertz/Mindlin contact
theory, assuming $E_{\mathrm{rubber}}/E_{\mathrm{PU}} \approx 1/10$ (e.g.\ Shore 60A
rubber vs.\ Shore 90A PU): contact radius $\chi \propto E^{-1/3}$, tangential
stiffness per unit load $\sigma_0 \propto G a/N \propto E^{2/3}$. Damping and
friction levels are indicative ranges only and should be replaced by measured
values.

\medskip
\begin{tabular}{@{}llll@{}}
\toprule
Constant & PU (base) & Rubber (estimate) & Basis \\
\midrule
$\mu$ & 0.3--0.8 (swept) & 0.6--0.9 & typical dry values \\
$r_s = \mu_s/\mu_c$ & 1.1 & 1.1--1.3 & indicative \\
$\hat\chi = \chi/h$ & 0.021 & 0.046 & $\chi\propto E^{-1/3}$ ($\times2.15$) \\
$\kappa_\chi = \chi/r_r$ & 0.345 & 0.74 & same \\
$\sigma_0$ [1/m] & 1640 & $\approx 350$ & $\sigma_0\propto E^{2/3}$ ($\times0.215$) \\
$\sqrt{\sigma_0 g}$ [rad/s] & 127 & 59 & \\
$\kappa_\sigma$ & 1.505 & 1.505 & assumed unchanged \\
$\zeta_b$ & 0.062 & $\approx$0.12--0.19 & higher loss factor ($\times$2--3) \\
$\Pi_b$ (relative) & 1 & $\approx 5$--6 & $\mu/\sigma_0$ \\
\bottomrule
\end{tabular}

\medskip
Qualitatively, rubber gives a longer pre-sliding range (larger $\Pi_b$, so a
wider linear zone), a slower and better-damped stick mode, and roughly twice the
spin coupling $\kappa_\chi$ (stronger Adamov spin--translation interaction).

\section*{Open notes}
\begin{itemize}
  \item Yaw moment arm uses $p_{y,i}$ while the contact velocity uses
    $p_{y,i} + DY_i$ (known; $\le 0.04\%$ of motor energy per
    \texttt{energy\_audit/README.md}).
  \item Roller torque equation: the $M_z$ term enters with sign $-M_z C_i$, whereas
    virtual-power consistency with $\omega_{z,i}$ (see the contact-kinematics
    section) requires $+M_z C_i$. To verify numerically.
\end{itemize}

\end{document}
